% Source coverage: physical PDF pp. 77--81 (printed pp. 54--58), from the
% Appendix B heading through the final Cartan-catalog paragraph before Problems.
% Displays: 23, all unnumbered. Footnotes/dividers: none. Citation: Ref. 26.
% Uncertainties: none after rendered-page review.

\chapterappendix{B}{The Cartan Catalog}

We present here without proof the complete catalog of simple Lie algebras,
worked out in its final form by E.~Cartan~\hyperref[ch15-ref-26]{[26]}. These will
be presented here in their `compact' form---that is, with generators that can be
faithfully represented by finite-dimensional Hermitian matrices. The Lie algebras
will be labelled with a subscript $n\geq1$ indicating their `rank'---the number of
independent commuting linear combinations of generators.

\noindent
$A_n$: This is the algebra of the special unitary group $SU(n+1)$, the group of
all unitary ($U^\dagger=U^{-1}$) unimodular ($\operatorname{Det}U=1$) matrices in
$n+1$ dimensions. Any such matrix that is infinitesimally close to the identity
may be expressed as
\[
 U=\mathbf{1}+iH ,
\]
with infinitesimal $H$ satisfying the conditions
\[
 H^\dagger=H,\qquad \operatorname{Tr}H=0 ,
\]
so $A_n$ is the Lie algebra of all Hermitian traceless matrices in $n+1$
dimensions. Any set of commuting Hermitian matrices may be simultaneously
diagonalized, and the maximum number of independent diagonal traceless matrices
in $n+1$ dimensions is $n$, so this is the rank of $A_n$. Any Hermitian matrix in
$n+1$ dimensions has $(n+1)^2$ independent real parameters ($n+1$ real numbers
on the main diagonal, and $n(n+1)/2$ complex numbers above the main diagonal,
equal to the complex conjugates of those below it), and one of these is eliminated
by the tracelessness condition, so the dimensionality of $A_n$ is
\[
 d(A_n)=(n+1)^2-1=n(n+2) .
\]
All of the $A_n$ are simple.

\noindent
$B_n$: This is the algebra of the unitary orthogonal group $O(2n+1)$, consisting
of all unitary ($\mathcal O^\dagger=\mathcal O^{-1}$) and orthogonal
($\mathcal O^{\mathsf T}=\mathcal O^{-1}$) and hence real matrices in $2n+1$ dimensions.
(The restriction to unitary matrices is sometimes indicated by calling the group
$UO(2n+1)$.) Any such matrix $\mathcal O$ that is infinitesimally close to the
identity may be expressed as
\[
 \mathcal O=\mathbf{1}+iA ,
\]
where $A$ is an infinitesimal matrix satisfying the conditions
\[
 A^*=-A=A^{\mathsf T} .
\]
(It would make no difference here if we restricted ourselves to the subgroup
$SO(2n+1)$, for which $\mathcal O$ is subject to the further condition
$\operatorname{Det}\mathcal O=1$, since any orthogonal matrix $\mathcal O$ that
is close to the identity would have $\operatorname{Det}\mathcal O=1$ anyway.)
Any set of commuting imaginary antisymmetric $(2n+1)$-dimensional matrices may
(by a common orthogonal transformation) be put in the form of a supermatrix
\[
 \begin{pmatrix}
  a_1\sigma_2&0&\cdots&0&0\\
  0&a_2\sigma_2&\cdots&0&0\\
  \vdots&\vdots&\ddots&\vdots&\vdots\\
  0&0&\cdots&a_n\sigma_2&0\\
  0&0&\cdots&0&0
 \end{pmatrix}
\]
where $\sigma_2$ is the usual $2\times2$ matrix
\[
 \sigma_2=\begin{pmatrix}0&-i\\i&0\end{pmatrix}
\]
and $a_1,\ldots,a_n$ are real. The rank of $B_n$ is thus evidently $n$. An
imaginary antisymmetric matrix is completely specified by the imaginary numbers
above the real diagonal, so its dimensionality is
\[
 d(B_n)=\frac{(2n+1)(2n)}{2}=n(2n+1) .
\]
All of the $B_n$ are simple.

There is an alternative definition of $O(N)$ that will help in understanding the
motivation for the next large set of simple Lie algebras. Instead of defining
$O(N)$ as the group of $N$-dimensional real matrices that satisfy the
orthogonality condition $\mathcal O^{\mathsf T}\mathcal O=\mathbf{1}$, it can just as well
be defined as the group of $N$-dimensional real matrices $\mathcal M$ that satisfy
the condition
\[
 \mathcal M^{\mathsf T}\mathcal P\mathcal M=\mathcal P ,
\]
where $\mathcal P$ is an arbitrary positive-definite real symmetric matrix. This
is because any such $\mathcal P$ may always be expressed as
$\mathcal P=\mathcal R^{\mathsf T}\mathcal R$ with $\mathcal R$ some real non-singular
matrix, so there is a similarity transformation that takes any matrix $\mathcal M$
satisfying the above condition into a real orthogonal matrix,
$\mathcal R\mathcal M\mathcal R^{-1}$. Thus we may let $\mathcal P$ be various
different real symmetric positive-definite matrices without changing the group.

\noindent
$C_n$: This is the algebra of the unitary symplectic group $USp(2n)$, the group
of unitary matrices $\mathcal M$ that leave an antisymmetric non-singular matrix
$\mathcal A$ invariant:
\[
 \mathcal M^{\mathsf T}\mathcal A\mathcal M=\mathcal A ,
\]
\[
 \mathcal A^{\mathsf T}=-\mathcal A,\qquad \operatorname{Det}\mathcal A\ne0 .
\]
(Note that in $d$ dimensions,
$\operatorname{Det}\mathcal A=\operatorname{Det}\mathcal A^{\mathsf T}
=(-)^d\operatorname{Det}\mathcal A$, so if $d$ is odd then
$\operatorname{Det}\mathcal A$ must vanish. Thus there is no $USp(d)$ unless
$d$ is even.) Any such antisymmetric non-singular (perhaps complex) matrix
$\mathcal A$ may be written in the standard form
\[
 \mathcal A=\mathcal R^{\mathsf T}\mathcal A_0\mathcal R ,
\]
where $\mathcal R$ is unitary and $\mathcal A_0$ is the supermatrix:
\[
 \mathcal A_0=\begin{pmatrix}0&\mathbf{1}\\-\mathbf{1}&0\end{pmatrix} .
\]
(Alternatively, we could take $\mathcal A_0$ as a block-diagonal supermatrix,
with $\sigma_1$s along the main diagonal.) Hence $USp(2n)$ may be described as
the group of unitary matrices $\mathcal S$ satisfying
\[
 \mathcal S^{\mathsf T}\mathcal A_0\mathcal S=\mathcal A_0
\]
since by a unitary transformation any such $\mathcal S$ can be transformed into
a unitary matrix $\mathcal M=\mathcal R^{-1}\mathcal S\mathcal R$ that satisfies
the previous condition $\mathcal M^{\mathsf T}\mathcal A\mathcal M=\mathcal A$. Any such
matrix $\mathcal S$ that is infinitesimally far from the identity may be written
\[
 \mathcal S=\mathbf{1}+i\mathcal H ,
\]
where $\mathcal H$ is an infinitesimal matrix satisfying the conditions
\[
 \mathcal H^\dagger=\mathcal H,\qquad
 \mathcal H^{\mathsf T}\mathcal A_0+\mathcal A_0\mathcal H=0 .
\]
The most general $2n$-dimensional matrix $\mathcal H$ satisfying these conditions
may be written as a supermatrix
\[
 \mathcal H=\begin{pmatrix}
  \mathcal A&\mathcal B\\
  \mathcal B^*&-\mathcal A^*
 \end{pmatrix} ,
\]
where the $n$-dimensional complex submatrices $\mathcal A$, $\mathcal B$ satisfy
\[
 \mathcal A^\dagger=\mathcal A,\qquad \mathcal B^{\mathsf T}=\mathcal B .
\]
A maximal set of commuting generators have $\mathcal A$ diagonal and
$\mathcal B$ zero, and so are of the form
\[
 \mathcal H=\operatorname{diag}(a_1,\ldots,a_n,-a_1,\ldots,-a_n)
\]
with $a_1,\ldots,a_n$ real. The rank of $C_n$ is thus evidently $n$. The
dimensionality of $C_n$ is the number $n^2$ of independent real parameters in the
Hermitian matrix $\mathcal A$, plus the number $2n(n+1)/2$ of independent real
parameters in the complex symmetric matrix $\mathcal B$
\[
 d(C_n)=n^2+2\cdot n(n+1)/2=n(2n+1) .
\]
All of the $C_n$ algebras are simple.

\noindent
$D_n$: This is the algebra of the unitary orthogonal group $O(2n)$, consisting of
all unitary orthogonal matrices in $2n$ dimensions. The discussion of $B_n$ can
be carried over to $D_n$, except that here any set of commuting generators can
be put in the form
\[
 \begin{pmatrix}
  a_1\sigma_2&0&\cdots&0\\
  0&a_2\sigma_2&\cdots&0\\
  \vdots&\vdots&\ddots&\vdots\\
  0&0&\cdots&a_n\sigma_2
 \end{pmatrix}
\]
so the rank is still $n$. Also, the dimensionality here is
\[
 d(D_n)=\frac{(2n)(2n-1)}{2}=n(2n-1) .
\]
All of the $D_n$ are simple except for $D_1$, which is the Abelian algebra with
just one generator, and $D_2$, which is the direct sum $B_1+B_1$.

\noindent
\textbf{Exceptional Lie Algebras:} In addition to the above classical Lie
algebras, there are just five special cases, the algebras $G_2(d=14)$;
$F_4(d=52)$; $E_6(d=78)$; $E_7(d=133)$; $E_8(d=248)$.

Not all of the classical Lie algebras are really different. There are just four
isomorphisms
\[
 A_1=B_1=C_1,\qquad C_2=B_2,\qquad A_3=D_3 .
\]
These correspond to isomorphisms among Lie groups of which these are the
algebras. However, isomorphisms like $B_1=A_1$, $B_2=C_2$, and $D_3=A_3$ do not
mean that $SO(3)$ is isomorphic to $SU(2)$, or that $SO(5)$ is isomorphic to
$USp(4)$, or that $SO(6)$ is isomorphic to $SU(4)$. Instead, $SU(2)$, $USp(4)$,
and $SU(4)$ are the simply connected covering groups of $SO(3)$, $SO(5)$, and
$SO(6)$. (Covering groups are discussed in Chapter 2.) Nevertheless, the
isomorphisms of the Lie algebras make it especially easy to construct the
double-valued fundamental spinor representations of $SO(3)$, $SO(5)$, and
$SO(6)$; they are just the defining representations of $SU(2)$, $USp(4)$, and
$SU(4)$, respectively. Also $SO(4)$ is isomorphic to $SO(3)\times SO(3)$, so its
double-valued spinor representation is just the defining representation of
$SU(2)\times SU(2)$. For $d\geq7$ the double-valued spinor representations of
$SO(d)$ must be constructed by other means. The simplest technique uses Clifford
algebras, discussed in Section 5.4.
